\chapter{Bitácora}\label{bitacora}

Reunión 27 agosto 2025

\begin{itemize}
    \item Revisión estructura trabajo/documento
    \item Revisión comentarios a propuesta de hitos
    \item Creación \href{https://notebooklm.google.com/notebook/c8490f2d-f084-4fe0-968a-d449fc4a9cde?authuser=1}{NotebookLM}
\end{itemize}

Reunión 3 de Septiembre 2025

\begin{itemize}
    \item Responder a lo comentado en los mails
    \item Posibilidad de que el cambio en el modelo del subastador realizado por Vicente sea tomar el c como variable y los resultados de su tesis sean agregados a mi modelo de subastador como parámetros. Posiblemente se tendrá que crear una restricción como que c este limitado a ser mayor o  igual a cero y menor o igual a un limite en el que haga sentido la inversión (no gastar más de lo que se va a obtener)
    \item \href{https://notebooklm.google.com/notebook/c8490f2d-f084-4fe0-968a-d449fc4a9cde?authuser=1}{NotebookLM}
\end{itemize}



escribir el problema amigo en marco teorico
abordar el problema del subastador a partir de lo poco definido de f
definir que será presicion en mi paper
definir que es el CAP para mi tesis

Reunión 24/09/25

\begin{itemize}
    \item Corregir orden marco teórico/ metodología
    \item Problema Amigo en marco teórico
    \item abordar que mejorar de la función del subastador basado en Amigo y porque. Definir bien presicion para mi tesis. Que será el CAP real también.
\end{itemize}

Reunión 01/10/25
\begin{itemize}
\item revisar Deep reserch de Gemini.
\item definir bien el espacio de Amigo donde se trabajara y sus implicancias
\item explicar A
\item grafico para modelo \citeA{munoz_lhuillier_inversion_2022}
\end{itemize}

Reunión 08/10/25
\begin{itemize}
\item no dejar de mencionar que se esta hablando del modelo Amigo Decir también que se asume que hay un mercado de productores que desean los permisos (pendiente). 
\item subir t=0 y t=1 explicación.
\item definir primero si se habla de secundario
\item interpretar bien la ecuación del subastador en bienestar económicamente
\item leer textos para mejorar introducción
\item subir bibliografías a Zotero 
\end{itemize}

Reunión 15/10/25
\begin{itemize}
\item Estudiar en profundidad punto 6.2 de Gabaix
\item Ver como es se aplica en Amigo
\item Escribir sobre las fallas del modelo de Vicente
\end{itemize}


Reunión 22/10/25
\begin{itemize}
\item Revisar críticamente el problema de la sección 6.2.2 de \citeA{gabaix_behavioral_2019} y escribirlo en función de $\theta$ y $CAP$
\item 
\item 
\end{itemize}

Reunión 29/10/25
\begin{itemize}
\item 
\item 
\item 
\end{itemize}

Reunión 5/11/25
\begin{itemize}
\item 
\item 
\item 
\end{itemize}
